\section{Opis problemu}
\label{sec:problem}

\subsection{Sformułowanie wymagań}
Rozważamy oświetlenie klatki schodowej sterowane z czterech przycisków
monostabilnych $S_1, S_2, S_3, S_4$ (po jednym na każdym piętrze) oraz
z czujnika ruchu oznaczonego dalej jako $M$. Wymagania
funkcjonalne są następujące:

\begin{enumerate}
    \item Światło można \textbf{włączyć} z dowolnego z czterech przycisków
          (gdy jest zgaszone).
    \item Jeżeli światło \textbf{świeci} i zostanie wciśnięty
          \emph{jakikolwiek} przycisk, światło gaśnie.
    \item Jeżeli światło świeci i przez \SI{2}{\minute} czujnik ruchu nie
          wykryje żadnego ruchu, światło gaśnie samoczynnie.
    \item Każde wykrycie ruchu przez czujnik $M$ \textbf{zeruje} odliczanie
          czasu (przedłuża świecenie o kolejne \SI{2}{\minute}).
\end{enumerate}

Punkty 1–2 sprowadzają się do klasycznego zachowania
\emph{przerzutnika typu T} (toggle): każde zbocze narastające sygnału
,,wciśnięto przycisk'' zmienia stan wyjścia $Q$. Punkt 3–4 realizuje
asynchroniczne zerowanie po upływie zadanego czasu bezruchu.

\subsection{Wielkości wejściowe i wyjściowe}
\begin{table}[H]
\centering
\caption{Zestawienie sygnałów układu.}
\begin{tabular}{cll}
\toprule
Symbol & Znaczenie & Typ \\
\midrule
$S_1,\dots,S_4$ & przyciski na piętrach & wejście, chwilowe \\
$M$              & czujnik ruchu     & wejście, poziom \\
$T$              & sygnał przekroczenia \SI{2}{\minute} bezruchu & wewnętrzny \\
$Q$              & stan oświetlenia (1 = świeci) & wyjście \\
\bottomrule
\end{tabular}
\end{table}

\subsection{Założenia projektowe}
\begin{itemize}
    \item Wszystkie przyciski są \emph{monostabilne}, tzn. sygnał jest
          aktywny tylko w trakcie naciśnięcia.
    \item Reakcja na ,,jakikolwiek przycisk'' modelowana jest bramką OR:
          $P = S_1 \lor S_2 \lor S_3 \lor S_4$.
    \item W części cyfrowej zakładamy, że sygnał $T$ jest już wytworzony
          przez timer (np. w mikrokontrolerze). W części 5 pokażemy go
          jako przekaźnik czasowy w schemacie drabinkowym.
    \item Stan wyjścia opisujemy zmienną $Q$; stan następny – $Q_{\text{next}}$.
\end{itemize}